\section{连续时间平稳分布与流量平衡}
\label{sec:06-Random-Processes/05-Continuous-Time-Markov-Chains-and-Queues/03}

两个状态的机器可以手解。把它换成一个有 20 个槽位的连接池——占用数从 0 到 20，请求到达占一个槽，处理完释放一个槽——你要问的仍然是同一类问题：长期看，各个占用水平各占多少时间？

按上一节的路子，先算 \(e^{tQ}\) 再令 \(t\to\infty\)。这条路走得通，但要对一个 \(21\times 21\) 的矩阵做指数运算，只为了拿一个不随时间变化的答案。即使退一步只解线性方程组 \(\pi Q=0\)，那也是 21 个方程。

其实一个方程都不用解。这一节要给的是一句可以在白板上口算的判断。

\subsection{平稳分布满足 \texorpdfstring{\(\pi Q=0\)}{piQ=0}}

行分布 \(\pi\) 若对所有 \(t\ge 0\) 满足 \(\pi P(t)=\pi\)，就称为平稳分布：初始就按 \(\pi\) 抽状态，此后任何时刻的分布都还是 \(\pi\)。在 \(t=0\) 处对两边求导，右端是常数导数为零，左端得到

\[
\pi Q=0,\qquad \sum_i\pi_i=1,\qquad \pi_i\ge 0.
\]

这是离散链 \(\pi P=\pi\) 的连续时间版本。差别在于离散链找的是特征值 1 的左特征向量，这里找的是特征值 0 的左特征向量——回想上一节的谱对应 \(\lambda\mapsto e^{t\lambda}\)，\(Q\) 的零特征值经指数化正好变成 \(P(t)\) 的特征值 1，两件事本来就是一件。

反过来的方向需要一点条件。有限状态链，或者更一般的保守、非爆炸链，满足 \(\pi Q=0\) 的概率向量确实给出平稳分布；对可能爆炸的无限状态生成矩阵，形式方程解出来了也未必对应一个真实的过程。

\subsection{读法是流入等于流出}

把 \(\pi Q=0\) 的第 \(j\) 个分量写开，利用 \(q_{jj}=-\sum_{k\ne j}q_{jk}\)：

\[
\sum_{i\ne j}\pi_i q_{ij}=\pi_j\sum_{k\ne j}q_{jk}.
\]

左边是长期单位时间内从别处流进 \(j\) 的概率流，右边是从 \(j\) 流出去的。平稳的含义就是每个状态的水位不再变化——进多少出多少。这叫全局平衡。

它比矩阵形式好记，也好用：写方程时不必先摆出 \(Q\)，只要在状态图上盯住一个节点，把指向它的箭头乘上源头的概率加起来，令其等于从它出发的箭头乘上自身概率之和。

注意全局平衡不要求每一对状态之间的流量都相等，只要求每个节点的总账平。允许存在环流：\(A\to B\to C\to A\) 方向上持续有净流量，而三个节点各自收支相抵。真实系统里这种环流很常见，比如只能按固定顺序推进的状态机。

若碰巧连每一对都平，即

\[
\pi_i q_{ij}=\pi_j q_{ji}\qquad\text{对所有 }i\ne j,
\]

就称满足细致平衡，链是可逆的——在平稳状态下把时间倒放，有限维分布不变。细致平衡蕴含全局平衡，反之不然，环流就是反例。

\subsection{生灭过程：一刀切下去就得到递推}

只允许在相邻状态间跳的链叫生灭过程：从 \(n\) 到 \(n+1\) 的出生率为 \(\lambda_n\)，从 \(n\) 到 \(n-1\) 的死亡率为 \(\mu_n\)。连接池、队列、冗余机组里的可用机器数，都是这个形状。

\begin{center}
\begin{tikzpicture}[x=1cm,y=1cm,font=\small,>=stealth]
  \foreach \i/\lab in {0/0,1/1,2/2,3/3,4/4}
    \draw[black!55] (1.7*\i,0) circle (0.34) node {\(\lab\)};
  \node[font=\footnotesize,black!60] at (7.6,0) {\(\cdots\)};
  \foreach \i in {0,1,2,3}
    \draw[->,blue!70] (1.7*\i+0.3,0.18) to[bend left=32] (1.7*\i+1.4,0.18);
  \foreach \i in {0,1,2,3}
    \draw[->,red!60] (1.7*\i+1.4,-0.18) to[bend left=32] (1.7*\i+0.3,-0.18);
  \draw[->,blue!70] (7.0,0.18) to[bend left=32] (7.4,0.28);
  \draw[->,red!60] (7.4,-0.28) to[bend left=32] (7.0,-0.18);
  \node[font=\footnotesize,blue!70] at (0.85,0.92) {\(\lambda_0\)};
  \node[font=\footnotesize,blue!70] at (2.55,0.92) {\(\lambda_1\)};
  \node[font=\footnotesize,blue!70] at (4.25,0.92) {\(\lambda_2\)};
  \node[font=\footnotesize,blue!70] at (5.95,0.92) {\(\lambda_3\)};
  \node[font=\footnotesize,red!60] at (0.85,-0.92) {\(\mu_1\)};
  \node[font=\footnotesize,red!60] at (2.55,-0.92) {\(\mu_2\)};
  \node[font=\footnotesize,red!60] at (4.25,-0.92) {\(\mu_3\)};
  \node[font=\footnotesize,red!60] at (5.95,-0.92) {\(\mu_4\)};
  \draw[black!70,dashed,thick] (4.25,-1.6) -- (4.25,-1.16);
  \draw[black!70,dashed,thick] (4.25,-0.7) -- (4.25,0.7);
  \draw[black!70,dashed,thick] (4.25,1.16) -- (4.25,1.6);
  \node[above,font=\footnotesize,black!75] at (4.25,1.6) {切口};
  \node[font=\footnotesize,black!80,align=center] at (4.25,-2.15) {跨过切口的只有一对箭头：\(\pi_2\lambda_2=\pi_3\mu_3\)};
\end{tikzpicture}

\emph{把状态分成左右两堆，只有一对箭头跨过切口。左堆的总概率既不增也不减，意味着两个方向的跨界流量必须相等——于是每一条边上都自动成立细致平衡。这不是额外假设，是``链没有环路''的直接后果：一维的链上，任何一次跨界都必须原路返回。}
\end{center}

切口论证给出的等式 \(\pi_n\lambda_n=\pi_{n+1}\mu_{n+1}\) 直接就是递推：

\[
\pi_{n+1}=\pi_n\frac{\lambda_n}{\mu_{n+1}},
\qquad\text{于是}\qquad
\pi_n=\pi_0\prod_{k=0}^{n-1}\frac{\lambda_k}{\mu_{k+1}}.
\]

20 个槽位的连接池，从 \(\pi_0\) 出发乘 20 次比值就完事，一个线性方程组都没解。最后一步是定 \(\pi_0\)，靠归一化 \(\sum_n\pi_n=1\)。

有限状态时归一化永远做得到。状态无限时它变成一道真正的检查：

\[
\sum_{n=0}^{\infty}\prod_{k=0}^{n-1}\frac{\lambda_k}{\mu_{k+1}}<\infty.
\]

级数发散就意味着不存在概率平稳分布——形式上的递推照样能写，但那串数加不出 1，系统不稳定。下一节的 M/M/1 里，这个级数是等比级数，收敛条件恰好是``到达率小于服务率''，稳定性判据就这么掉出来了。归一化不是收尾的手续，它是稳定性检查本身。

\subsection{回到那台机器}

两状态的故障修复模型是最短的生灭过程。切口在 0 与 1 之间，细致平衡给 \(\pi_0\lambda=\pi_1\mu\)，配上归一化：

\[
\pi_0=\frac{\mu}{\lambda+\mu},\qquad \pi_1=\frac{\lambda}{\lambda+\mu}.
\]

可用率等于修复速率占速率之和的比例。代入 \(\lambda=1/500\)、\(\mu=1/8\) 得 \(\pi_0\approx 98.4\%\)，与上一节瞬态解的极限 \(p(t)\to\lambda/(\lambda+\mu)=\pi_1\) 对上了——长时间之后初始状态被彻底忘掉，两条路殊途同归。

也可以顺手复核第一节那条换算。跳链在两个状态间硬性交替，平稳分布是 \(\widetilde\pi_0=\widetilde\pi_1=1/2\)；各除以离开速率再归一化，

\[
\frac{1/\lambda}{1/\lambda+1/\mu}=\frac{\mu}{\lambda+\mu}=\pi_0.
\]

跳链数的是访问次数，平稳分布数的是时间。拿跳链的结果当时间占比汇报，在这个例子里会把故障率报成 \(50\%\) 而非 \(1.6\%\)。

\subsection{解出来了不等于用得上}

有限、不可约的 CTMC 有唯一平稳分布，并且从任意初始分布收敛过去。这里没有离散链那种周期性的麻烦：随机的指数停留时间把节拍打散了，不会出现分布在几个子集间来回摆动、永不收敛的情形。这算是连续时间给的一点补偿。

无限状态则要多两道手续。链得是正常返的，形式解得能归一化，两者常常是同一件事的两种说法。可约的链还可能有多个闭类，每类一个平稳分布，它们的凸组合全都是平稳的——此时``唯一''失效，长期行为取决于初始状态落进了哪一类。

还有一件实践上的事：收敛与否是一回事，收敛得多快是另一回事。上一节已经看到，遗忘的时间尺度由 \(Q\) 的非零特征值决定。如果你的观测窗口短于这个尺度，报告平稳分布就是在报告一个系统还没到达的状态。稳定性不会因为你解出了一组非负数就自动成立。

下一节把生灭过程接到队列上。递推还是这个递推，但 \(\lambda_n\) 与 \(\mu_n\) 有了具体含义之后，归一化条件会变成一句关于服务器容量的硬话。
